% Conditional evidence-to-action reasoning for the hypothetical station case.
% Used in Section 5; standalone previews use sources/tikz/inferencechain.tex.
\begin{tikzpicture}[
  font=\small,
  node distance=5mm,
  >={Latex[length=2mm]},
  box/.style={draw, rounded corners=2pt, align=center,
    text width=7.25cm, minimum height=9mm, inner sep=4pt},
  side/.style={draw, rounded corners=2pt, align=center,
    text width=2.55cm, inner sep=4pt, font=\footnotesize},
  link/.style={->, thick},
  conditional/.style={->, thick, dashed},
  note/.style={font=\footnotesize, align=center, text width=7.3cm}
]
\node[box] (context) {\textbf{Assumed available station evidence}\\
  timetable, passenger or ticket flow,\\layout, service history};
\node[box, below=of context] (proxy) {\textbf{Proxy interpretation}\\
  which observations inform\\the use of a particular device?};
\node[box, below=of proxy] (load) {\textbf{Beliefs about usage, load and exposure}\\
  device-specific activity and accumulated burden};
\node[box, below=of load] (condition) {\textbf{Conditional beliefs about condition and events}\\
  exposure does not itself establish\\degradation or failure};
\node[box, below=of condition] (decision) {\textbf{Candidate maintenance actions}\\
  inspect earlier, seek better evidence,\\or retain the service rule};

\draw[link] (context) -- (proxy);
\draw[link] (proxy) -- (load);
\draw[conditional] (load) -- (condition);
\draw[link] (condition) -- (decision);

\node[side, left=5mm of load] (prior) {\textbf{Engineering priors and assumptions}\\
  traffic to usage;\\exposure to condition};
\draw[conditional] (prior.east) -- (load.west);
\draw[conditional] (prior.south) |- (condition.west);

\node[side, right=5mm of decision] (consequences) {\textbf{Consequences and constraints}\\
  service function,\\feasible actions,\\existing\\requirements};
\draw[link] (consequences.west) -- (decision.east);

\node[note, below=5mm of decision] {Uncertainty is retained through the reasoning chain.\\
  Dashed links require explicit engineering assumptions.};
\end{tikzpicture}
